% Seminário de Química nas Telecomunicações.
% Compilar com pdfLaTeX, duas passadas por causa do sumário, ou "make".
\documentclass[aspectratio=169,11pt]{beamer}

\usepackage[utf8]{inputenc}
\usepackage[T1]{fontenc}
% O microtype precisa de fonte escalável no pdfLaTeX.
\usepackage{lmodern}
\usepackage[brazil]{babel}
\usepackage{hyperref}

% Tema metropolis com paleta própria.

\usetheme[numbering=fraction,progressbar=frametitle]{metropolis}

\usepackage{graphicx}
\usepackage{booktabs}
\usepackage{array}
\usepackage{microtype}
\usepackage{xspace}
\usepackage{caption}
\usepackage{xcolor}

% Paleta: mAccent é a cor de contraste dos diagramas.
\definecolor{mPrimary}{HTML}{0A7E3B}
\definecolor{mPrimaryLight}{HTML}{7FCF88}
\definecolor{mPrimaryDark}{HTML}{064F26}
\definecolor{mNeutral}{HTML}{5B6770}
\definecolor{mAccent}{HTML}{B45309}

\setbeamercolor{structure}{fg=mPrimary}
\setbeamercolor{title separator}{fg=mPrimary}
\setbeamercolor{progress bar}{fg=mPrimary,bg=mPrimary!20}
\setbeamercolor{progress bar in head/foot}{fg=mPrimary,bg=mPrimary!20}
\setbeamercolor{progress bar in section page}{fg=mPrimary,bg=mPrimary!20}
\setbeamercolor{frametitle}{fg=white}
\setbeamercolor{alerted text}{fg=mPrimaryDark}
\setbeamercolor{example text}{fg=mPrimary}
\setbeamercolor{block title}{bg=mPrimary,fg=white}
\setbeamercolor{block body}{bg=white,fg=black}
\setbeamercolor{block title alerted}{bg=mPrimaryLight!45,fg=mPrimaryDark}
\setbeamercolor{block body alerted}{bg=white,fg=black}
\setbeamercolor{itemize item}{fg=mPrimary}
\setbeamercolor{itemize subitem}{fg=mPrimaryDark}
\setbeamercolor{itemize subsubitem}{fg=mNeutral}
\setbeamercolor{enumerate item}{fg=mPrimary}
\setbeamercolor{caption name}{fg=mPrimaryDark}

\metroset{
  sectionpage=progressbar,
  titleformat=smallcaps,
  block=fill,
  numbering=fraction
}

\usepackage{amsmath}
\usepackage{amssymb}
\usepackage{multirow}
\usepackage{chemfig}
\usetikzlibrary{arrows.meta,positioning,shapes.geometric,patterns,calc,decorations.pathreplacing}

% \qui{SiO_2} escreve a fórmula com o índice subscrito.
% Feito à mão para não depender do mhchem.
\newcommand{\qui}[1]{\ensuremath{\mathrm{#1}}}
\newcommand{\seta}{\ensuremath{\;\longrightarrow\;}}
\newcommand{\setaeq}{\ensuremath{\;\rightleftharpoons\;}}

% No lugar do siunitx.
\newcommand{\un}[2]{\ensuremath{#1\,\mathrm{#2}}}
\newcommand{\grau}[1]{\ensuremath{#1\,^{\circ}\mathrm{C}}}

\newcommand{\termo}[1]{\textcolor{mPrimaryDark}{\textbf{#1}}}

% Caixa que fecha um tópico ligando o assunto a telecom.
\newcommand{\gancho}[1]{%
  \begin{alertblock}{Por que isso importa em Telecomunicações?}#1\end{alertblock}}

\newcommand{\fonte}[1]{%
  \vfill{\scriptsize\textcolor{mNeutral}{Fonte: #1}}}

% \atomo{x}{y}{cor}{rótulo}: átomo dos diagramas de rede.
\newcommand{\atomo}[4]{%
  \node[circle,draw=#3,fill=#3!20,thick,minimum size=0.95cm,inner sep=0pt,
        font=\scriptsize\bfseries,text=#3] at (#1,#2) {#4};}


\title[Semicondutores em Telecom]{Materiais semicondutores na fabricação de dispositivos de telecomunicações}
\subtitle{Seminário de Química nas Telecomunicações}
\author[Victor Guerra, Alisson Pereira]{%
  \textbf{Alunos}: Victor Guerra, Alisson Pereira \\
  \textbf{Professor}: Maick Meneguzzo Prado}
\institute[IFSC]{\small IFSC, Câmpus São José \\ Engenharia de Telecomunicações}
\date{2 de setembro de 2026}

\begin{document}

\begin{frame}[plain]
  \titlepage
\end{frame}

\begin{frame}{Sumário}
  \footnotesize
  \begin{columns}[T]
    \begin{column}{0.5\linewidth}
      \tableofcontents[sections={1-4},hideallsubsections]
    \end{column}
    \begin{column}{0.5\linewidth}
      \tableofcontents[sections={5-8},hideallsubsections]
    \end{column}
  \end{columns}
\end{frame}

% Um arquivo por tópico. Victor nos ímpares, Alisson nos pares.
\section{Introdução}

\begin{frame}{Do quartzo da praia ao 5G}
  \begin{itemize}
    \item Todo enlace de telecomunicações, do smartphone ao transponder de
          satélite, depende de um punhado de \termo{cristais semicondutores}.
    \item O que separa a areia comum de um processador é \alert{química},
          mais precisamente estrutura cristalina, pureza e dopagem controlada.
    \item Um chip moderno concentra dezenas de bilhões de transistores em
          poucos \un{}{cm^2}, e nenhum deles funciona com uma impureza fora
          do lugar.
  \end{itemize}
  \vspace{0.3cm}
  \begin{block}{Roteiro do seminário}
    Material \seta{} estrutura \seta{} pureza \seta{} ligação \seta{} reação
    \seta{} dispositivo \seta{} preço.
  \end{block}
\end{frame}
                 % Abertura
\section{1. O que são esses materiais?}

\begin{frame}{Nem conduz, nem isola}
  \begin{columns}[T]
    \begin{column}{0.52\linewidth}
      \begin{itemize}
        \item No \termo{cobre} os elétrons estão soltos o tempo todo, então
              ele conduz sempre.
        \item No \termo{vidro} todos os elétrons estão presos nas ligações,
              então ele nunca conduz.
        \item No \termo{silício} os elétrons estão presos, mas soltam com
              pouca energia. Calor, luz ou uma tensão já bastam.
      \end{itemize}
      \vspace{0.2cm}
      \centering
      É o único dos três em que \alert{nós escolhemos o quanto ele conduz}.
    \end{column}
    \begin{column}{0.44\linewidth}
      \centering
      \begin{tikzpicture}[scale=0.68,transform shape]
        % cobre: as duas faixas se encostam, não há degrau
        \fill[mPrimary!30] (0,0) rectangle (1.6,1.1);
        \node[font=\scriptsize] at (0.8,0.55) {presos};
        \fill[mAccent!30] (0,1.1) rectangle (1.6,2.2);
        \node[font=\scriptsize] at (0.8,1.65) {soltos};
        \node[font=\scriptsize\bfseries] at (0.8,-0.45) {Cobre};
        \node[font=\tiny,text=mNeutral] at (0.8,-0.95) {sem degrau};

        \foreach \x/\gap/\nome/\tam in {3/1.1/Silício/pequeno, 6/2.6/Vidro/grande} {
          \fill[mPrimary!30] (\x,0) rectangle (\x+1.6,1.1);
          \node[font=\scriptsize] at (\x+0.8,0.55) {presos};
          \fill[mAccent!30] (\x,1.1+\gap) rectangle (\x+1.6,2.2+\gap);
          \node[font=\scriptsize] at (\x+0.8,1.65+\gap) {soltos};
          \draw[<->,thick,mNeutral] (\x+0.8,1.1) -- (\x+0.8,1.1+\gap);
          \node[font=\scriptsize\bfseries] at (\x+0.8,-0.45) {\nome};
          \node[font=\tiny,text=mNeutral] at (\x+0.8,-0.95) {degrau \tam};
        }
      \end{tikzpicture}
      \\[0.2cm]
      {\scriptsize O tamanho do degrau que o elétron precisa vencer é o
      que separa os três casos.}
    \end{column}
  \end{columns}
  \fonte{Callister, \emph{Ciência e Engenharia de Materiais}, cap.~18.}
\end{frame}

\begin{frame}{Sozinho o silício não serve}
  \begin{itemize}
    \item Em silício puro, só 1 átomo a cada 5 trilhões solta um elétron na
          temperatura ambiente. Conduz tão pouco que não dá para fazer nada
          com ele.
    \item A solução é sujar o cristal \alert{de propósito}, trocando um átomo
          a cada milhão por um vizinho da tabela periódica. Isso se chama
          \termo{dopagem}.
  \end{itemize}
  \vspace{0.2cm}
  \begin{columns}[T]
    \begin{column}{0.48\linewidth}
      \begin{block}{Fósforo, tipo N}
        \small Tem 5 elétrons para ligar e o silício só precisa de 4. O que
        sobra fica \termo{livre} para conduzir.
      \end{block}
    \end{column}
    \begin{column}{0.48\linewidth}
      \begin{block}{Boro, tipo P}
        \small Tem só 3 elétrons. A ligação que fica faltando é uma
        \termo{lacuna}, e ela também caminha pelo cristal.
      \end{block}
    \end{column}
  \end{columns}
  \vspace{0.2cm}
  \centering\small
  Uma impureza em um milhão faz a condutividade subir \alert{cem mil vezes}.
\end{frame}

\begin{frame}{Quem é quem na família}
  \begin{table}
    \centering
    \begin{tabular}{llc l}
      \toprule
      \textbf{Material} & & \textbf{Degrau (eV)} & \textbf{Onde aparece} \\
      \midrule
      Silício            & Si   & 1,1 & Chips, memórias, quase tudo \\
      Arseneto de gálio  & GaAs & 1,4 & Amplificadores de celular, LED \\
      Nitreto de gálio   & GaN  & 3,4 & Antenas 5G, carregadores \\
      Fosfeto de índio   & InP  & 1,3 & Lasers de fibra óptica \\
      \bottomrule
    \end{tabular}
  \end{table}
  \begin{itemize}\small
    \item Degrau pequeno solta elétron fácil, mas vaza corrente e esquenta.
    \item Degrau grande aguenta tensão alta, e por isso o GaN vai para as
          antenas 5G.
  \end{itemize}
  \gancho{As duas bases são \termo{silício e gálio}. O silício domina por ser
  barato; o gálio entra onde é preciso alta frequência ou luz.}
\end{frame}
                   % 1, Victor
% Os slides deste tópico entram abaixo do \section.
\section{2. Estrutura química}
   % 2, Alisson
\section{3. Pureza dos materiais}

\begin{frame}{A pureza que o silício precisa ter}
  \begin{columns}[T]
    \begin{column}{0.52\linewidth}
      \begin{itemize}
        \item O silício da eletrônica tem \alert{99,999999999\%} de pureza,
              os \termo{onze noves}.
        \item Isso é menos de \termo{um átomo estranho por bilhão}.
        \item É um grão de açúcar dissolvido em uma piscina olímpica.
      \end{itemize}
    \end{column}
    \begin{column}{0.44\linewidth}
      \small
      \begin{table}
        \centering
        \begin{tabular}{ll}
          \toprule
          \textbf{Uso} & \textbf{Pureza} \\
          \midrule
          Liga metálica & 98\% \\
          Painel solar  & 99,9999\% \\
          Eletrônica    & 11 noves \\
          \bottomrule
        \end{tabular}
      \end{table}
    \end{column}
  \end{columns}
  \vspace{0.4cm}
  \centering\small
  A dopagem usa um átomo em um milhão. Se a sujeira estiver nessa mesma
  proporção, ela toma o lugar do dopante.
\end{frame}

\begin{frame}{Como o silício é purificado}
  \begin{enumerate}
    \item \termo{Forno}: areia mais carbono a \grau{2000}. Sai um silício com
          98\% de pureza.
    \item \termo{Destilação}: o silício vira gás, o gás é destilado e depois
          volta a ser sólido, já limpo.
    \item \termo{Cristalização}: o cristal é puxado devagar do silício
          derretido, e a sujeira fica no líquido.
  \end{enumerate}
  \vspace{0.4cm}
  \gancho{No amplificador de uma antena, impureza vira ruído, e mais ruído
  significa menos alcance para a estação.}
\end{frame}
                         % 3, Victor
% Os slides deste tópico entram abaixo do \section.
% O enunciado pede a junção PN (boro tipo P, fósforo tipo N).
\section{4. Ligações químicas}
   % 4, Alisson
\section{5. Reações químicas}

\begin{frame}{Obtenção do silício a partir da areia}
  \begin{block}{Forno a arco elétrico, \grau{\sim 2000}}
    \[ \qui{SiO_2}\,(s) + 2\,\qui{C}\,(s) \seta \qui{Si}\,(l) + 2\,\qui{CO}\,(g) \]
  \end{block}
  \vspace{0.2cm}
  \begin{itemize}
    \item O carbono tira o oxigênio da areia, numa reação de
          \termo{oxirredução}.
    \item Gasta muita energia, e por isso a produção fica onde a energia é
          barata.
    \item O silício sai com 98\% de pureza, ainda \alert{muito longe} do que
          a eletrônica precisa.
  \end{itemize}
\end{frame}

\begin{frame}{Purificação pelo processo Siemens}
  \begin{block}{1. O silício vira gás, a \grau{300}}
    \[ \qui{Si}\,(s) + 3\,\qui{HCl}\,(g) \seta \qui{SiHCl_3}\,(g) + \qui{H_2}\,(g) \]
  \end{block}
  \begin{block}{2. O gás volta a ser silício, a \grau{1100}}
    \[ \qui{SiHCl_3}\,(g) + \qui{H_2}\,(g) \seta \qui{Si}\,(s) + 3\,\qui{HCl}\,(g) \]
  \end{block}
  \vspace{0.2cm}
  \centering\small
  Purificar sólido é difícil, mas o gás pode ser \termo{destilado} como se
  destila álcool. A segunda reação é a primeira \alert{ao contrário}, e o
  silício volta sem a sujeira.
\end{frame}

\begin{frame}{Reações usadas na fabricação}
  \begin{block}{Cria a camada isolante sobre o silício}
    \[ \qui{Si}\,(s) + \qui{O_2}\,(g) \seta \qui{SiO_2}\,(s) \]
  \end{block}
  \vspace{0.2cm}
  \begin{block}{Forma o cristal de arseneto de gálio}
    \[ \qui{Ga(CH_3)_3}\,(g) + \qui{AsH_3}\,(g) \seta \qui{GaAs}\,(s) + 3\,\qui{CH_4}\,(g) \]
  \end{block}
  \vspace{0.3cm}
  \centering\small
  Para desenhar o circuito, o ácido fluorídrico abre janelas na camada
  isolante sem atacar o silício.
\end{frame}
     % 5, Victor
% Os slides deste tópico entram abaixo do \section.
% O enunciado pede smartphones, satélites e roteadores.
\section{6. Aplicabilidade}
         % 6, Alisson
\section{7. Valores comerciais}

\begin{frame}{O preço da pureza}
  \small
  \begin{table}
    \centering
    \begin{tabular}{lrl}
      \toprule
      \textbf{Produto} & \textbf{Preço (US\$)} & \textbf{Observação} \\
      \midrule
      Quartzo, matéria-prima        & 0,03 a 0,10 / kg  & abundante \\
      Silício metalúrgico (2N)      & 1,5 a 2,5 / kg    & commodity \\
      Polisilício eletrônico (11N)  & 25 a 60 / kg      & varia com o ciclo \\
      Wafer de Si polido 300\,mm    & 90 a 150 / unid.  & virgem, sem processo \\
      Wafer de GaAs 150\,mm         & 250 a 450 / unid. & \\
      Wafer de SiC 150\,mm          & 600 a 1.200 / unid. & gargalo de oferta \\
      \bottomrule
    \end{tabular}
    \caption{Ordens de grandeza, referência de 2024 e 2025.}
  \end{table}
  \centering
  Do quartzo ao polisilício eletrônico o valor sobe cerca de
  \alert{$1000\times$}, e o que se comprou foi \termo{pureza}, não massa.
\end{frame}

\begin{frame}{Onde o valor se concentra}
  \begin{columns}[T]
    \begin{column}{0.52\linewidth}
      \begin{itemize}\small
        \item Um wafer de \un{300}{mm} custa cerca de US\$\,120 cru; processado
              em nó avançado sai da fábrica por \alert{US\$\,15 a 20 mil}.
        \item Dele saem centenas de chips, então o custo unitário cai, mas a
              máscara de um nó de 3\,nm passa de US\$\,20 milhões.
        \item Uma fábrica de ponta custa de US\$\,15 a 20 bilhões.
      \end{itemize}
    \end{column}
    \begin{column}{0.44\linewidth}
      \begin{block}{Mercado global (2024)}
        \small
        \begin{itemize}
          \item Semicondutores: US\$\,630 bi
          \item Fatia de comunicações: 30\%
          \item GaN para RF: US\$\,1,5 bi
        \end{itemize}
      \end{block}
    \end{column}
  \end{columns}
  \vspace{0.1cm}
  \gancho{O gálio é subproduto do refino de \termo{bauxita} e mais de 90\% da
  oferta mundial vem da China, e o preço saltou de US\$\,300 para além de
  US\$\,600/kg após os controles de exportação de 2023.}
\end{frame}
 % 7, Victor

\section*{Referências}

\begin{frame}[shrink=15]{Referências}
  % Fontes menores para caber em um slide só
  \setbeamerfont{bibliography item}{size=\scriptsize}
  \setbeamerfont{bibliography entry author}{size=\scriptsize}
  \setbeamerfont{bibliography entry title}{size=\scriptsize}
  \setbeamerfont{bibliography entry location}{size=\scriptsize}
  \setbeamerfont{bibliography entry note}{size=\scriptsize}
  \begin{thebibliography}{9}
    \setbeamertemplate{bibliography item}[text]

    \bibitem{callister}
    CALLISTER, W. D.; RETHWISCH, D. G.
    \newblock \emph{Ciência e Engenharia de Materiais: uma Introdução}.
    \newblock 10.~ed. Rio de Janeiro: LTC, 2021. Cap.~18.

    \bibitem{atkins}
    ATKINS, P.; JONES, L.; LAVERMAN, L.
    \newblock \emph{Princípios de Química: questionando a vida moderna}.
    \newblock 7.~ed. Porto Alegre: Bookman, 2018.

    \bibitem{sze}
    SZE, S. M.; LI, Y.; NG, K. K.
    \newblock \emph{Physics of Semiconductor Devices}.
    \newblock 4th ed. Hoboken: Wiley, 2021.

    \bibitem{streetman}
    STREETMAN, B. G.; BANERJEE, S. K.
    \newblock \emph{Solid State Electronic Devices}.
    \newblock 7th ed. Harlow: Pearson, 2016.

    \bibitem{shriver}
    WELLER, M. \emph{et al.}
    \newblock \emph{Química Inorgânica}.
    \newblock 7.~ed. Porto Alegre: Bookman, 2020. Cap.~13--14.

    \bibitem{usgs}
    U.S. GEOLOGICAL SURVEY.
    \newblock \emph{Mineral Commodity Summaries: Gallium; Silicon}.
    \newblock Reston: USGS, 2025.

    \bibitem{wsts}
    WORLD SEMICONDUCTOR TRADE STATISTICS.
    \newblock \emph{WSTS Semiconductor Market Forecast}. 2025.
    \newblock Disponível em: \url{https://www.wsts.org}.
  \end{thebibliography}
\end{frame}


\section*{Conclusão}

\begin{frame}{Fechando o ciclo}
  \begin{itemize}
    \item O semicondutor é um cristal comum cuja condutividade nós aprendemos
          a \termo{controlar}, e quem define esse controle é a banda proibida.
    \item A \termo{pureza de 11N} é o que permite que a dopagem intencional,
          na casa de ppm, signifique alguma coisa dentro do cristal.
    \item Uma sequência de reações bem conhecidas leva do quartzo ao wafer,
          passando por redução, destilação do \qui{SiHCl_3}, cristalização e
          deposição.
    \item O preço acompanha a \termo{pureza e o controle de processo}, e não a
          quantidade de material, e por isso o valor sobe mil vezes ao longo
          do caminho.
  \end{itemize}
\end{frame}

\begin{frame}[plain]
  \vfill
  \centering
  {\Large Obrigado! \\[0.3cm] Perguntas?}
  \vfill
\end{frame}


\end{document}
